\documentclass[../main.tex]{subfiles}
\graphicspath{{\subfix{../Images/}}}

\begin{document}
\onecolumn
% =================================================================
% Sections V-VI -- Method and Implementation (skeleton)
%
% Companion to problem_formulation.tex / problem_formulation_fixed.tex.
% Realises every forward reference that Section IV already makes into this
% file's namespace: sec:method, sec:method:storage, sec:method:controllers,
% sec:method:rl, sec:method:model, sec:method:resolution, sec:method:training,
% sec:implementation, sec:impl:instances, sec:impl:logging, sec:impl:validation
% all match exactly, so \ref/\eqref calls already written in Section IV
% resolve against this file with no changes needed there.
%
% CONTRACT (from Method_and_Implementation.md, restated for enforcement here)
%   Section V may only refer to Section IV by equation number, never restate a
%   definition. If a paragraph below needs to define something Section IV
%   already defines (a task, a conflict, the coupling condition, the storage
%   capacity constraint, Pi, F), that is a bug in this file -- it should be an
%   \eqref, not a restatement. Every equation label already used by
%   problem_formulation_fixed.tex is treated as read-only here.
%
% This is a SKELETON: every numeric value in the parameter tables is marked
% TBD and is meant to be filled in after the first stable runs (see
% Method_and_Implementation.md, "Recommended order of work", Pass 5). The
% algorithmic choices (which greedy heuristic, which priority rule, which
% reward shape) are concrete enough to implement and test, but are Method-level
% choices, not yet validated by any experiment -- treat them as the first
% reasonable draft, not as settled design.
%
% REQUIRES BIBLIOGRAPHY ENTRIES (natbib \citet/\citep form; see the equivalent
% note in problem_formulation_fixed.tex if main.tex does not load natbib):
%   knippenberg2021, ng1999policy, devlin2012dynamic -- already required by
%     problem_formulation_fixed.tex and GAPS.tex; add once, used in all three.
%   @article{schulman2017ppo,
%     author = {Schulman, John and Wolski, Filip and Dhariwal, Prafulla and
%               Radford, Alec and Klimov, Oleg},
%     title  = {Proximal Policy Optimization Algorithms},
%     journal = {arXiv preprint arXiv:1707.06347}, year = {2017} }
%   @misc{rllibmultiagent,
%     author = {{Ray Project}}, title = {Multi-Agent Environments --- RLlib},
%     howpublished = {\url{https://docs.ray.io/en/latest/rllib/multi-agent-envs.html}},
%     year = {2026} }
%   @misc{pettingzooparallel,
%     author = {{Farama Foundation}}, title = {Parallel API --- PettingZoo},
%     howpublished = {\url{https://pettingzoo.farama.org/api/parallel/}},
%     year = {2026} }
%   @inproceedings{sartoretti2019primal,
%     author = {Sartoretti, Guillaume and Kerr, Justin and Shi, Yunfei and
%               Wagner, Glenn and Satish Kumar, T. K. and Koenig, Sven and
%               Choset, Howie},
%     title  = {{PRIMAL}: Pathfinding via Reinforcement and Imitation
%               Multi-Agent Learning},
%     journal = {IEEE Robotics and Automation Letters}, volume = {4}, number = {3},
%     pages = {2378--2385}, year = {2019} }
%
% Labels defined by this file:
%   Section: sec:method, sec:method:overview, sec:method:storage,
%     sec:method:controllers, sec:method:rl, sec:method:model,
%     sec:method:resolution, sec:method:training, sec:implementation,
%     sec:impl:simulator, sec:impl:instances, sec:impl:logging,
%     sec:impl:cost, sec:impl:validation
%   Equation: eq:storagegreedy eq:assignmentrule eq:features eq:mask
%     eq:reward eq:msgpass eq:readout
%   Figure: fig:methodloop fig:architecture
%   Table: tab:storageparams tab:rlparams tab:modelparams tab:trainparams
% =================================================================

\section{Method}
\label{sec:method}

Section~\ref{sec:problem} fixed what the objects are: the sets, the
constraints, the coupling condition, and the objective
$\mathcal{J}(F,\pi)$ as a framing device. This section realises every
component that section left open --- the storage rule $F \in \mathcal{F}$
(\ref{sec:pf:storage}), the three controllers $\pi \in \Pi$
(\ref{sec:pf:controllers}), the reward $R_i$
(\eqref{eq:objective}), and the conflict-resolution mechanism whose
existence \eqref{eq:vertexconflict}--\eqref{eq:swapconflict} already assume
(\ref{sec:pf:agents}). Nothing below changes a definition from
Section~\ref{sec:problem}; every reference to an object defined there is by
equation or section number.

% -----------------------------------------------------------------
\subsection{Overview}
\label{sec:method:overview}

All three controllers of Table~\ref{tab:information} are embedded in the same
six-stage loop, so that a controller swap changes only stages 2--3 and the
comparison stays matched. At each timestep: (1)~orders drawn from
$P_{\mathrm{ord}}$ are converted into new tasks \eqref{eq:task} and enter
$\mathcal{Q}_t$; (2)~task assignment (Section~\ref{sec:method:controllers},
identical across all three architectures) matches free agents to tasks in
$\mathcal{Q}_t$; (3)~routing produces one intended action
$u_i(t)\in\mathcal{U}(\ell_i(t))$ per agent, using whichever information the
active architecture permits (Table~\ref{tab:information}); (4)~the
conflict-resolution operator (Section~\ref{sec:method:resolution}) revises
the joint intended action into one satisfying
\eqref{eq:vertexconflict}--\eqref{eq:swapconflict}, flagging any agent it
overrides; (5)~the environment applies the resolved joint action, advancing
$\ell_i(t)$, the lifecycle sets \eqref{eq:lifecycle}, and the log
(Section~\ref{sec:impl:logging}); (6)~if $t$ is a storage epoch, the storage
rule additionally applies \eqref{eq:storageupdate}. Figure~\ref{fig:methodloop}
shows the loop; a reader who stops here already knows the shape of the system
that the rest of this section fills in piece by piece.

\begin{figure}[t]
  \centering
  \fbox{\parbox{0.9\linewidth}{\centering\vspace{2cm}
    \emph{[Figure: the six-stage loop above, drawn as a cycle diagram in the
    style of Figure~\ref{fig:loop}, with the three controller-specific
    boxes (centralised / section-based / decentralised) shown as
    interchangeable inserts at stage 3.]}
    \vspace{2cm}}}
  \caption{The environment step. Stages 1, 2, 4, 5 and 6 are shared by every
  controller; only stage 3 (routing) differs, which is what makes a result
  difference attributable to the routing architecture.}
  \label{fig:methodloop}
\end{figure}

% -----------------------------------------------------------------
\subsection{Storage update rule}
\label{sec:method:storage}

Three named rules realise $F \in \mathcal{F}$ (\ref{sec:pf:storage}), ordered
as a ladder rather than presented as one sophisticated rule, because the gap
between the second and third isolates exactly this thesis's contribution: the
value of \emph{traffic} information over \emph{demand} information alone.

\textbf{$F_{\mathrm{fix}}$} is the identity rule already fixed in
\eqref{eq:storageupdate}'s surrounding text; it is not re-derived here.

\textbf{$F_{\mathrm{dem}}$} (demand-only) is the classical slotting
heuristic: sort $\mathcal{K}$ by $\hat\rho_t(k)$ descending, sort
$V_{\mathrm{str}}$ by $\min_{g\in V_{\mathrm{del}}} d_G(v,g)$ ascending, and
place SKUs into vertices in that joint order, skipping a vertex once it is at
capacity \eqref{eq:feasiblestorage}. This is $\arg\min_{x\in\mathcal{X}}
D(x;\hat\rho_t)$ from \eqref{eq:storageobjective} with $\beta=\chi=0$, solved
greedily rather than exactly.

\textbf{$F_{\mathrm{cng}}$} (congestion-aware) extends $F_{\mathrm{dem}}$'s
ranking with a congestion penalty: SKU $k$'s score at candidate vertex $v$
becomes
\begin{equation}
  \mathrm{score}(k,v) = \hat\rho_t(k)\cdot d_G\bigl(v, \textstyle\min_{g}\bigr)
  \;+\; \beta \sum_{e \in \mathrm{path}(v)} \hat\mu_t(e)\,\hat w_t(e) ,
  \label{eq:storagegreedy}
\end{equation}
where $\mathrm{path}(v)$ is the shortest path used to reach $v$ from the
nearest delivery point and $\beta \ge 0$ weights congestion exposure against
access distance; SKUs are placed in ascending score order, subject to the same
capacity check. \eqref{eq:storagegreedy} realises $D+\beta K$ from
\eqref{eq:storageobjective} greedily.

The reassignment term $R$ is realised as a \emph{hard cap} rather than a
smooth penalty: at most $\nu$ units may be relocated relative to
$x_{t-\Delta}$ per epoch, i.e.\ $R(x,x_{t-\Delta}) = 0$ if the number of
relocated units is $\le \nu$ and $R = \infty$ otherwise, so any finite
$\chi > 0$ in \eqref{eq:storageobjective} enforces the cap and $\chi$ itself
is fixed rather than swept. When $F_{\mathrm{dem}}$ or $F_{\mathrm{cng}}$
would relocate more than $\nu$ units, only the $\nu$
highest-marginal-value relocations are applied, processed in strictly
descending marginal score, ties broken by SKU id then vertex id --- this is
the tie-breaking rule that \eqref{eq:storageupdate} requires to be a genuine
function rather than a set-valued $\arg\min$ (see \ref{sec:pf:storage}'s
requirement). $\Delta=\infty$ is not a separate rule to implement:
$F_{\mathrm{fix}}$ already covers it exactly.

\begin{table}[h]
  \centering
  \caption{Storage-rule parameters. TBD columns are filled after the first
  stable runs (Method\_and\_Implementation.md, Pass 5).}
  \label{tab:storageparams}
  \begin{tabular}{lll}
    \toprule
    Parameter & Meaning & Value \\
    \midrule
    $\beta$ & congestion weight in \eqref{eq:storagegreedy} & TBD (sweep) \\
    $\nu$ & max relocated units per epoch & TBD (sweep) \\
    $\Delta$ & storage epoch length & TBD (sweep, per \ref{sec:pf:coupling}) \\
    \bottomrule
  \end{tabular}
\end{table}

% -----------------------------------------------------------------
\subsection{Controller realisations}
\label{sec:method:controllers}

The three architectures of Table~\ref{tab:information} realise
$\pi=(\pi^{\text{assign}},\pi^{\text{route}}) \in \Pi$
(\ref{sec:pf:controllers}). Task assignment is realised once and shared by
all three, so that a result difference is never confounded with assignment
quality: at every timestep, free agents are matched to tasks in
$\mathcal{Q}_t$ ordered by $r_j$ ascending (oldest first), each matched
greedily to the free agent minimising $d_G(\ell_i(t), s_j)$,
\begin{equation}
  i^\star(j) = \arg\min_{i \,:\, a_i \text{ free}} d_G(\ell_i(t), s_j) ,
  \label{eq:assignmentrule}
\end{equation}
removing $a_{i^\star(j)}$ from the free pool before considering the next
task. This is a greedy heuristic for bipartite matching, not the Hungarian
algorithm's optimum; using the optimum instead is a one-line ablation, not a
change to $\Pi$.

\textbf{Centralised} $\pi^{\text{route}}$: prioritised planning
(space-time A$^\star$) over the full graph, agents planned in a fixed order
(earliest $y_j$ first); each agent after the first plans around the
space-time cells already reserved by higher-priority agents for the current
planning horizon.

\textbf{Section-based} $\pi^{\text{route}}$: the same prioritised planner run
independently inside each zone $Y_q$ of \eqref{eq:partition}, using only
zone-local state. The boundary protocol: an agent whose planned route crosses
$Y_q \to Y_{q'}$ yields at the boundary to whichever zone currently has fewer
active agents, ties broken by agent id --- one explicit, stated rule, not left
implicit.

\textbf{Decentralised} $\pi^{\text{route}}$: the learned policy
$\pi_\theta(\cdot\mid o_i(t))$ of Sections~\ref{sec:method:rl}
and~\ref{sec:method:model}.

% -----------------------------------------------------------------
\subsection{Reinforcement-learning formulation}
\label{sec:method:rl}

Four components realise the decentralised regime's $\pi_\theta$: how
$o_i(t)$ \eqref{eq:observation} is encoded, the action space and its mask,
the reward, and the objective (already \eqref{eq:objective}, not restated).

\textbf{Observation features.} Each visible vertex $v \in
\mathcal{G}_i^{(d)}(t)$ contributes a feature vector
\begin{equation}
  f(v,t) = \bigl(\, \mathrm{onehot}(\mathrm{role}(v)),\;\; \eta_i(v,t),\;\;
                    \mathbb{1}[v \text{ occupied by another agent}] \,\bigr) ,
  \label{eq:features}
\end{equation}
where $\mathrm{role}(v) \in \{V_{\mathrm{str}}, V_{\mathrm{del}}, V_{\mathrm{ep}},
\text{transit}\}$ (\ref{sec:pf:env}) and $\eta_i(v,t)$ is
\eqref{eq:potential}. Each message from $a_j$ with $(a_i,a_j)\in
E^{\mathcal{A}}_t$ \eqref{eq:commgraph} carries $(d_G(\ell_i(t),\ell_j(t)),\,
\eta_j(\ell_j(t),t))$. The congestion feature $\delta_i(t)$
\eqref{eq:congestion} is appended once, agent-level rather than per-node.

\textbf{Action space and masking.} With $d_{\max} = \max_{v\in V_{\mathrm{mov}}}
|N^{+}(v)|$, the policy head outputs $d_{\max}+1$ logits (one wait, $d_{\max}$
move slots). At a state with $\ell_i(t)=v$, logits for slots beyond
$|N^{+}(v)|$ --- i.e.\ for $\mathcal{U}_i \setminus \mathcal{U}(v)$
(\ref{sec:pf:game}) --- are set to $-\infty$ before the softmax,
\begin{equation}
  \Pr[u_i(t) = \cdot] = \mathrm{softmax}\bigl(\, \ell(\cdot) + M(v,\cdot) \,\bigr),
  \qquad M(v,\cdot) = \begin{cases} 0 & \cdot \in \mathcal{U}(v) \\
                                     -\infty & \cdot \in \mathcal{U}_i\setminus\mathcal{U}(v) \end{cases}
  \label{eq:mask}
\end{equation}
which is what makes each individual action legal \emph{by construction}
rather than in expectation. RLlib's newer multi-agent API does not ship
action masking as a built-in and it must be implemented as a custom module or
connector \citep{rllibmultiagent}; budget implementation time for it
separately from the training-procedure work of
Section~\ref{sec:method:training}.

\textbf{Reward.} Writing $\Phi_i(t) = -\eta_i(\ell_i(t),t)$ for the
progress potential built from \eqref{eq:potential},
\begin{equation}
  R_i(t) = \gamma\Phi_i(t+1) - \Phi_i(t)
           \;+\; r_{\mathrm{deliver}}\,\mathbb{1}[a_i \text{ completes a task at } t]
           \;-\; r_{\mathrm{blk}}\,\mathbb{1}[a_i \text{ is overridden at } t]
           \;-\; r_{\mathrm{cng}}\,\delta_i(t) .
  \label{eq:reward}
\end{equation}
The first term is potential-based shaping in the form of
\citet{ng1999policy}: it does not alter the underlying stochastic game's
equilibria, and the guarantee extends to dynamic potentials
\citep{devlin2012dynamic} -- required here specifically, since $q_i(t)$ and
hence $\Phi_i$ change at every reassignment \eqref{eq:lifecycle}, making
$\Phi_i$ non-stationary by construction. A naive ``reward for decreasing
distance'' term lacks this guarantee and admits policies that loiter near a
target to farm shaping reward; \eqref{eq:reward} does not, by the same
argument. The override indicator $\mathbb{1}[a_i \text{ is overridden at }
t]$ is produced by the conflict-resolution operator of
Section~\ref{sec:method:resolution} and is the same flag that operationally
defines $\omega_j(t)$ in \eqref{eq:waiting} --- one piece of instrumentation
serves both the reward and the run-level metric, logged once
(Section~\ref{sec:impl:logging}). Do not finalise $r_{\mathrm{blk}}$ before
the resolution operator is verified: the penalty is meaningless until
overrides are detected correctly, and a mis-scaled feasibility penalty is
among the most common causes of a MAPF policy that learns to stand still.
The last term, $-r_{\mathrm{cng}}\,\delta_i(t)$, is a plain per-step penalty,
\emph{not} shaping --- unlike the first term, it can in principle shift the
equilibrium, which is an accepted, stated trade-off for giving the policy a
direct congestion-avoidance signal, not an oversight.

\begin{table}[h]
  \centering
  \caption{RL formulation parameters. TBD columns are filled after the first
  stable runs.}
  \label{tab:rlparams}
  \begin{tabular}{lll}
    \toprule
    Parameter & Meaning & Value \\
    \midrule
    $r_{\mathrm{deliver}}, r_{\mathrm{blk}}, r_{\mathrm{cng}}$ & reward weights, \eqref{eq:reward} & TBD (sweep) \\
    $\gamma$ & discount, \eqref{eq:objective} & TBD \\
    $d$ & observation depth, \eqref{eq:localsubgraph} & TBD (sweep) \\
    $r_{\mathrm{com}}$ & communication radius, \eqref{eq:commgraph} & TBD \\
    $r_{\mathrm{cng}}$ & congestion radius, \eqref{eq:congestion} & TBD \\
    \bottomrule
  \end{tabular}
\end{table}

% -----------------------------------------------------------------
\subsection{Model architecture}
\label{sec:method:model}

The overall policy/value architecture descends from the same PRIMAL lineage
\citep{sartoretti2019primal} that \citet{knippenberg2021} already adapted from
a CNN-on-lattice encoder to a single-layer GCN-on-arbitrary-graph encoder; the
step taken here is the same kind of substitution one level further, from a
single GCN pass to explicit multi-round message passing so that information
can travel more than one hop within $\mathcal{G}_i^{(d)}(t)$ per decision.
The encoder embeds $\mathcal{G}_i^{(d)}(t)$ into $z_i$ via $Q$ rounds of
message passing over node features \eqref{eq:features}: writing
$h_v^{(0)} = f(v,t)$,
\begin{align}
  m_v^{(q)} &= \phi_q\bigl(h_v^{(q-1)},\, \{h_w^{(q-1)} : w \in
               \mathcal{G}_i^{(d)}(t),\, (w,v)\in E \text{ or } (v,w)\in E\}\bigr),
  \label{eq:msgpass} \\
  h_v^{(q)} &= \psi_q\bigl(h_v^{(q-1)},\, m_v^{(q)}\bigr), \qquad q = 1,\dots,Q,
  \notag \\
  z_i &= h_{\ell_i(t)}^{(Q)} \;\Vert\; \mathrm{AGG}\bigl(\{\text{messages
        from } a_j : (a_i,a_j)\in E^{\mathcal{A}}_t\}\bigr) \;\Vert\; \delta_i(t) ,
  \label{eq:readout}
\end{align}
concatenating the agent's own final-layer embedding, an aggregation (e.g.\
mean) of incoming messages, and the congestion feature. $\pi_\theta(\cdot\mid
o_i(t))$ and a value head $V_\theta(o_i(t))$ (for training,
Section~\ref{sec:method:training}) are fully connected layers with ReLU
activations over $z_i$, masked per \eqref{eq:mask} before the softmax.
$Q=0$ collapses \eqref{eq:msgpass}--\eqref{eq:readout} to raw node/message
features with no propagation at all, which ablates \emph{both} local
structural aggregation and inter-agent communication at once; if the two are
meant to be isolated separately (structure vs.\ communication), a cleaner
ablation replaces the single $Q$ with two independent round-counts, one for
propagation within $\mathcal{G}_i^{(d)}(t)$ and one for propagation across
$E^{\mathcal{A}}_t$ --- noted here as an open refinement, not yet adopted, to
avoid overstating what a single $Q=0$ run actually isolates.

\begin{figure}[t]
  \centering
  \fbox{\parbox{0.9\linewidth}{\centering\vspace{2cm}
    \emph{[Figure: local observation $\mathcal{G}_i^{(d)}(t)$ and incoming
    messages $\to$ $Q$ rounds of $\phi_q,\psi_q$ $\to$ $z_i$ $\to$ FC+ReLU
    $\to$ masked softmax / value head, in the style of the source paper's
    Figure~3 but with the message-passing block expanded.]}
    \vspace{2cm}}}
  \caption{Per-agent model. Weights are shared across all agents (parameter
  sharing, \ref{sec:pf:controllers}).}
  \label{fig:architecture}
\end{figure}

\begin{table}[h]
  \centering
  \caption{Model architecture parameters.}
  \label{tab:modelparams}
  \begin{tabular}{lll}
    \toprule
    Parameter & Meaning & Value \\
    \midrule
    $Q$ & message-passing rounds, \eqref{eq:msgpass} & TBD (sweep, incl.\ $Q=0$) \\
    hidden width & $\phi_q,\psi_q$, FC layers & TBD \\
    \bottomrule
  \end{tabular}
\end{table}

% -----------------------------------------------------------------
\subsection{Conflict resolution operator}
\label{sec:method:resolution}

Masking \eqref{eq:mask} makes each individual action legal; it does not make
the \emph{joint} action legal, since two masked-legal individual choices can
still violate \eqref{eq:vertexconflict} or \eqref{eq:swapconflict}. A
deterministic operator removes the remainder. At the start of each episode, a
priority permutation $\sigma$ of $\{1,\dots,m\}$ is drawn once from the
episode seed and held fixed for the episode. At each timestep, given the
proposed joint action $(u_1(t),\dots,u_m(t))$:
\begin{enumerate}
  \item process agents in order $\sigma$;
  \item for agent $a_{\sigma(k)}$, tentatively accept its proposed successor
    vertex if it violates neither \eqref{eq:vertexconflict} nor
    \eqref{eq:swapconflict} against the already-finalised successors of
    $a_{\sigma(1)},\dots,a_{\sigma(k-1)}$;
  \item otherwise, override $a_{\sigma(k)}$'s action to $\mathrm{wait}$
    (always legal, \eqref{eq:actions}) and set its override flag for $t$.
\end{enumerate}
Because waiting never conflicts with an already-finalised successor other
than by occupying the same current vertex (impossible, since the joint state
at $t$ was itself collision-free), this always terminates in a feasible joint
action, realising the guarantee stated in \ref{sec:pf:agents} that a
collision-free joint action always exists. The operator is a fixed function
of the state and the episode seed, with no additional randomness --- this is
exactly the determinism that Section~\ref{sec:pf:scope} states is required
for $P$ to be stochastic only through order arrivals; that scope claim holds
because of this construction, not independently of it. Determinism here is
also what makes paired-seed comparisons across $(F,\pi)$ combinations valid:
the same seed and the same realised trace up to any policy difference imply
the same resolution decisions.

The override flag set in step~3 is logged (Section~\ref{sec:impl:logging})
and used identically in two places already committed to above: as
$\mathbb{1}[a_i\text{ is overridden at }t]$ in \eqref{eq:reward}, and to
compute $\omega_j(t)$ in \eqref{eq:waiting}.

% -----------------------------------------------------------------
\subsection{Training procedure}
\label{sec:method:training}

Training uses PPO \citep{schulman2017ppo} with full parameter sharing: every
agent maps to one shared policy id, RLlib's standard multi-agent scaling
lever \citep{rllibmultiagent}. Seeds are split disjointly between training
and evaluation before any run starts. Two training regimes are distinguished
and reported as an experimental variable, not folded together: training under
$F_{\mathrm{fix}}$ only (then evaluated against all three storage rules), and
training under the same storage rule used at evaluation (``matched''
training). The latter is harder: training under an adaptive $F$ makes the
task distribution non-stationary from the learner's perspective
(\ref{sec:pf:coupling}), since \eqref{eq:coupling} routes pickups differently
as $x_t$ changes underneath the policy. $c_{\mathrm{wait}}$
(\ref{sec:pf:env}, footnote) is reported alongside training configuration,
since it is a parameter of the study, not a training hyperparameter.

\begin{table}[h]
  \centering
  \caption{PPO training parameters.}
  \label{tab:trainparams}
  \begin{tabular}{lll}
    \toprule
    Parameter & Meaning & Value \\
    \midrule
    learning rate, clip $\epsilon$, GAE $\lambda$ & PPO core & TBD \\
    minibatch size, epochs/update, rollout length & PPO core & TBD \\
    entropy coeff., value-loss coeff. & PPO core & TBD \\
    train/eval seed counts & disjoint split & TBD \\
    \bottomrule
  \end{tabular}
\end{table}

% =================================================================
\section{Implementation}
\label{sec:implementation}

This section describes how the model of Sections~\ref{sec:problem}
and~\ref{sec:method} is built and verified. It contains no new modelling
decisions --- anything that would count as one belongs in
Section~\ref{sec:method} instead, per the contract stated at the top of this
file.

% -----------------------------------------------------------------
\subsection{Simulator}
\label{sec:impl:simulator}

All-pairs shortest paths are precomputed once per instance, since $G$
(\ref{sec:pf:env}) is static and $d_G$ is queried every timestep by both the
observation builder \eqref{eq:potential} and the storage rule
\eqref{eq:storagegreedy}. Task generation, assignment, routing, resolution,
storage update, and logging are separate modules with explicit interfaces, so
a controller swap (Section~\ref{sec:method:controllers}) touches only the
routing module. Agents act simultaneously; observations, rewards and actions
are returned as dictionaries keyed by agent id, matching the POSG formalism
of \eqref{eq:posg} directly and the conventions of RLlib's
\texttt{MultiAgentEnv} and PettingZoo's Parallel API
\citep{rllibmultiagent,pettingzooparallel}.

% -----------------------------------------------------------------
\subsection{Instance generation}
\label{sec:impl:instances}

Warehouse graphs are generated from parameters (aisle count, aisle length,
cross-aisle count, one-way fraction) rather than hand-drawn, which is what
substantiates the ``non-lattice'' claim of \ref{sec:pf:env}. Each generated
instance is checked against the finiteness/connectivity assumption of
\ref{sec:pf:env} (the choice between the weaker per-pair assumption and full
strong connectivity, noted as open there, is resolved concretely here: a
generated instance is discarded if any vertex needed as a target is
unreachable from any occupiable vertex) and, only when a baseline requires
it, against the well-formedness condition of \ref{sec:pf:env} ($\ge m$
non-task endpoints, path between any two not through a third). The discard
rate is reported. The order generator ($P_{\mathrm{ord}}$,
\ref{sec:pf:tasks}) is seeded separately from the policy and the environment,
so the same order stream replays identically under every controller and
storage rule, which is what makes the paired comparison in
\eqref{eq:rqformal} meaningful.

% -----------------------------------------------------------------
\subsection{Instrumentation}
\label{sec:impl:logging}

Every metric is accumulated online, not recomputed at the horizon. Per-edge
counters give $\mu_T(e)$ and hence $H_T$ \eqref{eq:entropy}. A per-task record
of $r_j$, $y_j$, arrival-at-$s_j$ time, and $d_j$ gives $\zeta_j$
\eqref{eq:servicetime}, with the waiting and active phases reported
separately per \eqref{eq:lifecycle}. $\omega_j(t)$ \eqref{eq:waiting} is
logged directly from the conflict-resolution operator's override flag
(\ref{sec:method:resolution}): for the agent serving each active task, one bit
per timestep records whether it chose $\mathrm{wait}$ or was overridden,
summed per task to give $W_T$ with no separate instrumentation needed. Backlog
$B_T$ is logged every timestep, not only at the horizon, so a diverging queue
is visible in the trace rather than only in the final summary.

% -----------------------------------------------------------------
\subsection{Computational cost}
\label{sec:impl:cost}

Runtime is a dependent variable, not an aside: wall-clock per decision is
logged separately for assignment \eqref{eq:assignmentrule} and routing, on
the same hardware across controllers. For the learned controller, training
time and inference time are reported separately, since the practical case for
a learned router rests specifically on cheap inference, not on cheap training.

% -----------------------------------------------------------------
\subsection{Correctness and reproducibility}
\label{sec:impl:validation}

Four assertions run on \emph{every} experiment, not once in development:
\begin{enumerate}
  \item no vertex or swap conflict on the realised trace
    (\eqref{eq:vertexconflict}--\eqref{eq:swapconflict});
  \item every realised storage configuration lies in $\mathcal{X}$
    \eqref{eq:feasiblestorage};
  \item every released task is a member of exactly one of $\mathcal{Q}_t$,
    $\mathcal{B}_t$, $\mathcal{C}_t$ at every timestep --- which
    \eqref{eq:lifecycle} now makes a checkable invariant rather than merely
    an assumed one, since disjointness follows from the interval definition
    and can be asserted directly against $y_j$ and $d_j$;
  \item bitwise replay determinism under fixed seeds, following from
    Section~\ref{sec:method:resolution}'s determinism and
    Section~\ref{sec:impl:instances}'s separated seeding.
\end{enumerate}
These four checks are cheap, and they are also the answer to ``how do you
know the simulator implements the model of Section~\ref{sec:problem}.''
Section~\ref{sec:experiments} reports the resulting estimates of
\eqref{eq:rqformal}, having verified through this section that the estimates
are of the model actually specified, not of an implementation drift from it.

\end{document}
